\chapter{Improper Error Handling}
\label{ch:error_handling}

\begin{concept}{Error Handling}{error-handling-def}
    Error handling is a mechanism used to "resolve" (handle) errors that arise during the execution of a program. It deals with error events to avoid program crashes. Crucially, error handling is a part of the overall security of an application.
\end{concept}

Improper error handling occurs when the software does not handle some error conditions, leaving the program in an unsecure state. This unsecure state might eventually produce a crash, resulting in a potential \gls{dos}. Furthermore, unhandled errors might be accompanied by the disclosure of sensitive information about the code itself when inappropriate error messages propagate to the end user.

\section{Information Disclosure and Reconnaissance}
An attack always begins with a Reconnaissance phase in which the attacker will try to gather as much technical information about the target. Improper error handling and unhandled errors can help attackers in collecting useful information.

For example, verbose exception reports might disclose the technology stack in use (such as Opensymphony, Struts2, and Tomcat versions) via stack traces rendered to the user. Similarly, the disclosure of a \gls{sql} query error can expose the site installation path, which can be used to identify an injection point. 

\section{Mishandling Exceptional Conditions in Code}

Failing to properly capture and manage errors leads to fail-open behaviors and unpredicted application states. Below are common programming pitfalls.

\subsection{Unchecked Return Values (C)}
When opening a file using \texttt{fopen}, the function will return a \texttt{NULL} pointer if the operation fails. 
\begin{lstlisting}[language=C]
FILE* f = fopen(szFilename, "r");
// read data from f
fclose(f);
\end{lstlisting}
If \texttt{szFilename} is invalid (e.g., manipulated by an attacker), \texttt{f} will be \texttt{NULL}. If the program proceeds to use this null pointer for file operations, the function will crash. To fix this, the pointer must be validated (\texttt{if (f == NULL)}), and the error should be reported to the caller, where it is supposed to be handled properly.

\subsection{Misunderstanding \glsfmtshort{api} Error Mechanisms (C++)}
Developers frequently mishandle exceptions by misunderstanding the underlying library's documentation. 

\textbf{Masked Exceptions:}
\begin{lstlisting}[language=C++]
BigThing* p = new (std::nothrow) BigThing[14999];
\end{lstlisting}
The \texttt{std::nothrow} flag indicates that \texttt{new} shall not throw an exception on failure, but will return a null pointer instead. Consequently, if allocation fails, any surrounding \texttt{catch(std::bad\_alloc\& e)} branch won't be executed. The subsequent use of the null pointer \texttt{p} will cause a crash.

\textbf{Catching the Wrong Exception:}
When allocating a \texttt{CString} object in \gls{mfc} programs, an insufficient memory error throws a \texttt{CMemoryException}. If the developer writes a catch block anticipating \texttt{std::bad\_alloc\&}, the exception will bypass the handler entirely.

\subsection{Useless Error Checks (C)}
Developers sometimes check the wrong variable for failure indicators.
\begin{lstlisting}[language=C]
char* p = strncpy(dest, szBuf, 19);
if (p) { ... }
\end{lstlisting}
The developer wrongly believes that the return value of \texttt{strncpy} is \texttt{NULL} in case of an error. However, \texttt{strncpy} returns a pointer to the start of the destination string regardless of the outcome. Therefore, the check \texttt{if (p)} is useless for error handling.

\section{Mitigation Strategies}
To ensure robust security posture during error events, developers must:
\begin{itemize}
    \item Handle errors and exceptions in the code in appropriate ways.
    \item Never mask exceptions that may corrupt the program's state.
    \item Use assertions in the code to document assumptions and/or unexpected conditions or for situations that should never occur.
    \item Check function return values for errors.
    \item Filter error messages provided to users.
\end{itemize}

\section{Real-World Impact (\glsfmtshortpl{cve})}
Improper error handling leads to significant real-world vulnerabilities:
\begin{itemize}
    \item \textbf{CVE-2018-1000083:} Ajenti version 2 contains an Improper Error Handling vulnerability in the Login \gls{json} request. By sending a malformed \gls{json}, the tool responds with a traceback error that leaks a path of the server.
    \item \textbf{CVE-2018-14907:} The Web server in 3CX version 15.5.8801.3 is vulnerable to Information Leakage because of improper error handling in stack traces, demonstrating the discovery of a full pathname.
    \item \textbf{CVE-2024-36481:} In the Linux kernel, a function might return \texttt{NULL} or an error via a macro, but its caller only checked for the \texttt{NULL} condition, missing the explicit error return.
    \item \textbf{CVE-2018-1999036:} An exposure of sensitive information vulnerability exists in Jenkins \gls{ssh} Agent Plugin. It exposes the \gls{ssh} private key password to users with permission to read the build log.
\end{itemize}